\documentclass[11pt, reqno]{article}

\usepackage{amsmath, amsthm, amssymb}
\usepackage{enumitem}
\usepackage{tcolorbox}
\usepackage{hyperref}
\usepackage{tikz}
\usepackage{tikz-cd}
\usepackage{pgfplots}
\pgfplotsset{compat=1.18}
\usetikzlibrary{arrows.meta}
\usepackage{mathrsfs}
\usepackage{fancyhdr}
\usepackage[bottom=0.75in, top=1in, left=0.5in, right=0.5in]{geometry}
\usepackage{array}   % for \newcolumntype macro
\newcolumntype{L}{>{$}l<{$}}

\theoremstyle{plain}
\newtheorem*{theorem}{Theorem}
\newtheorem*{proposition}{Proposition}
\newtheorem{exercise}{Exercise}
\newtheorem*{lemma}{Lemma}
\newtheorem*{corollary}{Corollary}

\theoremstyle{definition}
\newtheorem*{definition}{Definition}
\newtheorem*{example}{Example}

\theoremstyle{remark}
\newtheorem*{remark}{Remark}

\renewcommand{\phi}{\varphi}
\renewcommand{\epsilon}{\varepsilon}
\renewcommand{\emptyset}{\varnothing}

\newcommand{\RR}{\mathbb{R}}
\newcommand{\ZZ}{\mathbb{Z}}
\newcommand{\NN}{\mathbb{N}}
\newcommand{\CC}{\mathbb{C}}
\newcommand{\QQ}{\mathbb{Q}}

\DeclareMathOperator{\ima}{\text{im}}

\begin{document}

\topmargin=-40pt
\rhead{Henry Woodburn}
\lhead{Math 656}
\renewcommand{\headrulewidth}{1pt}
\renewcommand{\headsep}{20pt}
\thispagestyle{fancy}

{\Huge \bfseries \noindent Homework 3}

\begin{enumerate}
    \item[3.] We can modify the proof of (i) to remove the factor of $2$. For any $0 < \delta < 1$,
    we can choose $y_n$ in $Y_n$ such that 
    \[  
        |y_n| = 1, \quad |y_n - y| < \delta
    \]
    for all $y \in Y_{n-1}$. This turns (25) into 
    \[
        |\textbf{C} y_n - \textbf{C} y_m| \geq \epsilon |\lambda_n|,
    \]

    and (26) becomes 
    \[
        N(\epsilon) \leq C(\epsilon, \delta^{-1}\textbf{C}(B)).
    \]
    Since this holds for all $\epsilon < 1$, we can remove the factor of epsilon altogether.

    \item[4.] Since $(\zeta - \textbf{C})$ has a pole of order $i$ at $\lambda$, the function 
    $(\zeta - \lambda)^i (\zeta - \textbf{C})^{-1}$ is analytic in a neighborhood of $\lambda$. This means 
    we can write a power series 
    \[
        (\zeta - \lambda)^i (\zeta - \textbf{C})^{-1} = \sum_{n = 0}^\infty A_n (\zeta - \lambda)^n,
    \]
    and thus 
    \[
        (\zeta - \textbf{C})^{-1} = \sum_{n = 0}^infty A_n(\zeta - \lambda)^{n-i},
    \]
    or 
    \[
        (\zeta - \textbf{C})^{-1} = \sum_{n = -i}^\infty \tilde{A}_n (\zeta - \lambda)^n.
    \]

    From theorem 6 of chapter 17, we see that by how the operators $\tilde(A)_n$ are constucted, by a line
    integral surrounding an isolated point in the spectrum, we see that these operators are projections, by the result
    concerning the functional calculus. I could not finish.

    \item[6.] We know from functional calculus that the spectrum of $\textbf{B}^n$ is exactly the $n$-th power 
    of the spectrum of $\textbf{B}$. We apply a previous result to see that the number of eigenvalues of $\textbf{B}^n$
    greater than $\epsilon^n$, $N'(\epsilon^n)$, can be estimated by
    \[
        N'(\epsilon^n) \leq C(\epsilon, \textbf{B}^n(B)),
    \]
    and we know the set of eigenvalues of $\textbf{B}^n$ of norm greater than $\epsilon^n$ is the same 
    as the set of eigenvalues of $\textbf{B}$ with norm greater than $\epsilon$. Then we have
    \[
        N(\epsilon) \leq C(\epsilon^n, \textbf{B}^n(B)).
    \]

    \item[8.] Consider the hilbert space $\ell^2(\mathbb{N}) = \ell^2$, and let $M_i$ be the projection onto 
    the first $i$ coordinates. $M_i$ is compact since the range is finite dimensional. Moreover, $M_i$ converges
    strongly (pointwise for every $x \in \ell^2)$ to the identity operator on $\ell^2$. However, the identity operator
    is clearly not compact. 

    \item[9.] Let $X$ be a Banach space. Suppose $M_n$ converges to $M$ strongly and $C$ is a compact operator. Then
    $M_n x$ converges in norm to $Mx$ for all $x \in X$. We want to show 
    \[
        \lim_{n \to \infty}\sup_{x \in B}|CM_n x - CM x| = \lim_{n \to \infty} \sup_{x \in B}|C(M_n - M)x| = 0
    \]

    and

    \[
        \lim_{n \to \infty}\sup_{x \in B}|(M_n - M)Cx| = 0.
    \]
    Since the image of the unit ball under $C(M_n - M)$ and $(M_n - M)C$ is compact, this must be true. Suppose not,
    that the limit does not go to zero. Then we can choose a sequence $x_k$ and a subsequence $n_k$ such that
    \[
        |C(M_{n_k} - M)x_k| > \epsilon
    \]
    for all $k$, and an analagous expression for the opposite composition. But since the image is compact, there must 
    be a convergent subsequence $k_j$. Since the norms of $M_n$ are bounded, it follows that $\lim_j |C(M_{n_{k_j}} - M)x_{k_j}| = 0$,
    contradicting the earlier lower bound. Then $CM_n$ must converge to $CM$ uniformly, and $M_n C$ must converge
    to $MC$ uniformly. 

\end{enumerate}

\end{document}